\documentclass[paper=a4, fontsize=13pt, twoside]{scrartcl} % A4 paper and 11pt font size
\linespread{1.5}
%\renewcommand{\familydefault}{\sfdefault}
\usepackage{indentfirst}
\usepackage[T1]{fontenc} 
\usepackage{times}
\usepackage[polish]{babel} 
\usepackage[utf8]{inputenc}
\usepackage{amsmath,amsfonts,amsthm} 
%\usepackage{color}
\usepackage[pdftex]{graphicx}
%\usepackage{pdfpages}
\usepackage{graphicx}
\usepackage{subcaption}
\usepackage[draft,nosingleletter]{impnattypo}
\usepackage{sectsty} % Allows customizing section commands
\usepackage{setspace}
\allsectionsfont{\normalfont \scshape} % Make all sections centered, the default font and small caps

\usepackage{longtable}

\usepackage[left=30mm,right=25mm,top=25mm,bottom=25mm,bindingoffset=5mm]{geometry}

\usepackage{caption}
\usepackage{listings}

\usepackage{lipsum}

\usepackage{fancyhdr}
\pagestyle{fancy}
\fancyhead[RO, LE]{\thepage}
\fancyhead[RE, LO]{}
\fancyfoot{}
\renewcommand\headrulewidth{0pt}
\renewcommand\footrulewidth{0pt}

\fancypagestyle{firststyle}{
\fancyhead[RO, LE]{}
\fancyhead[RE, LO]{}
\renewcommand{\footrulewidth}{1pt}
\cfoot{\textrm{\textbf{\data}}}
}

\fancypagestyle{taskstyle}{
\fancyhead[RO, LE]{}
\fancyhead[RE, LO]{}
\renewcommand{\footrulewidth}{0pt}
}

\renewcommand\headrulewidth{0pt}
\renewcommand\footrulewidth{0pt}

\setlength\parindent{27pt}

\newcommand{\TODO}[1]{
	{\large \color{red}{\underline{\textbf{TODO: #1}}}} 
}

\newcommand{\inserttitlepage}{
    \thispagestyle{firststyle}
    \begin{center}
    \textrm{\LARGE{\textbf{WOJSKOWA AKADEMIA TECHNICZNA}}\\
    \normalsize\textbf{im. Jarosława Dąbrowskiego}
    \rule[4mm]{\linewidth}{1pt}
    \LARGE{\textbf{WYDZIAŁ CYBERNETYKI}}}
    \end{center}
    
    \vspace{5mm}
    
    \begin{center}
        \includegraphics[width=40mm]{./defs/WAT.png}
    \end{center}

	\vspace{5mm}

    \begin{center}
        \textbf{{\fontsize{36}{10} \selectfont \textrm{PRACA DYPLOMOWA}}}\\
        \vspace{5mm}
        \textbf{\textrm{\LARGE{\studia}}}
    \end{center}
    
    \begin{longtable*}{p{.25\textwidth} p{.75\textwidth}}
        \textrm{\Large{Temat pracy:}} & \textrm{\textbf{\Large{\temat}}} \\
    \end{longtable*}
    
    \begin{singlespace}
    \begin{center}
        \textrm{\textbf{\Large{\kierunek}}}\\
        \textrm{...............................................................................................................\\{{\scriptsize{(kierunek studiów)}}}}\\~\\
        \textrm{\textbf{\Large{\specjalnosc}}}
        \textrm{...............................................................................................................\\\scriptsize{(specjalność)}}\\
    \end{center}
    \end{singlespace}
    % \vspace{30mm}
    \begin{longtable*}{p{.5\textwidth} p{.5\textwidth}}
    \textrm{Dyplomant:} & 	\textrm{Promotor pracy:}\\
    \textrm{\textbf{\autor}}  & \textrm{\textbf{\promotor}} \\
    \ifdefined\konsultant
    ~ & ~\\ 
    ~ & \textrm{Konsultant:} \\
    ~ & \textrm{\textbf{\konsultant}}\\
    \fi
    \end{longtable*}
    
    \clearpage
    
    %\newpage
  
    %\large{\textbf{Oświadczenie}}\\
    
    %\normalsize{"Wyrażam zgodę na udostępnianie mojej pracy przez Archiwum  WAT"}
    %~\\
    
    %\begin{longtable*}{p{.6\textwidth} c}
    %Data ...................... & 	....................................\\
    %~ & (podpis)\\
   
    %~ & ~\\ 
   
    %~ & Pracę przyjąłem \\
    
    %~ & ~\\
    %~ & ~\\
        
    %~ & Promotor Pracy\\
    %~ & \promotor \\
    %\end{longtable*}
	
    \includegraphics[width=1.0\textwidth ,height=1.0\textheight , keepaspectratio]{./files/os-pod.pdf}

    \includegraphics[width=1.0\textwidth ,height=1.0\textheight , keepaspectratio]{./files/zad1.pdf}
  
    \includegraphics[width=1.0\textwidth ,height=1.0\textheight , keepaspectratio]{./files/zad2.pdf}
   
    %\includegraphics[width=\pagewidth ,height=\pageheight, keepaspectratio]{./files/zad2.pdf}
     
    %\large{Zadania do pracy dyplomowej:}
    
    %\normalsize{\zadania}
    
    %\newpage
    
    %\textit{\normalsize{Ta strona zostanie zamieniona na oryginał zadań do pracy dyplomowej}}
    
    \newpage
   

    \renewcommand\refname{Bibliografia}
}

\newcommand{\rysunek}[4]{
    \begin{figure}[h]
        \begin{center}
            \includegraphics[width=#2\textwidth]{./files/#1}
            \caption{#3 \\ \textit{Źródło: #4}}
            \label{#1}
        \end{center}
    \end{figure}   
}

\newcommand{\podrysunek}[3]{
    \begin{subfigure}[t]{.4\textwidth}
    	\centering
        \includegraphics[width=\linewidth]{./files/#1}
    \end{subfigure}   
}
\newcommand{\wykres}[4]{
    \begin{figure}[h]
        \begin{center}
            \includegraphics[width=#2\textwidth]{./files/#1}
            \caption{#3. \textit{Źródło: #4}}
            \label{#1}
        \end{center}
    \end{figure}   
}

\captionsetup[table]{name=Tabela}
\newcommand{\tabela}[5]{                            
        \begin{longtable}{#4}
        \caption{#2.}\\
            #5
        \end{longtable}
}

\captionsetup[lstlisting]{margin=0pt}

\renewcommand{\lstlistingname}{Kod}
\newcommand{\kod}[3]{                            
       \lstinputlisting[caption=#2. \textit{Źródło: #3}\label{#1}, basicstyle=\ttfamily]{./files/#1}
}
